\emph{Setting.} By Fact \claimref{standard-model.imported-one-loop-running-trace-ratio}(a) each $a_{i}$ is linear in $t = \ln(\mu/M_{Z})/2\pi$ with slope $-b_{i}$ below every threshold, and content $\Dl^{(k)}$ switched on at $t_{k} \ge 0$ adds $\Dl^{(k)}_{i}$ to the coefficient above $t_{k}$; integrating gives the displayed $a_{i}(t) = a_{i} - b_{i}t - \sum_{k}\Dl^{(k)}_{i}(t - t_{k})_{+}$. The record values are those of Proposition \claimref{standard-model.spectrum-run}: $a_{2} = 0.2312 \cdot 128 = 18496/625$, $a_{1} = \tfrac35(128 - a_{2}) = 184512/3125$, $a_{3} = 500/59$. With $b_{1} - b_{3} = \tfrac{41}{10} + 7 = \tfrac{111}{10}$ these give exactly
\[
t_{13} \;=\; \frac{a_{1} - a_{3}}{b_{1} - b_{3}} \;=\; \frac{18647416}{4093125} \;=\; 4.5558\ldots, \qquad
a_{3} + 7t_{13} \;=\; \frac{165219412}{4093125} \;=\; 40.365\ldots,
\]
which is $40.4$ to the precision stated.

\emph{(i).} Put $d(t) := a_{1}(t) - a_{3}(t)$. Since $a_{1} - a_{3} = (b_{1} - b_{3})t_{13}$,
\[
d(t) \;=\; (b_{1} - b_{3})(t_{13} - t) \;-\; \sum_{k}\bigl(\Dl^{(k)}_{1} - \Dl^{(k)}_{3}\bigr)(t - t_{k})_{+}.
\]
Every term of the sum is $\ge 0$ by the hypothesis $\Dl^{(k)}_{1} \ge \Dl^{(k)}_{3}$, so $d(t) \le (b_{1} - b_{3})(t_{13} - t)$. The function $d$ is continuous and piecewise linear with slope $-(b_{1} - b_{3}) - \sum_{t_{k} < t}(\Dl^{(k)}_{1} - \Dl^{(k)}_{3}) \le -\tfrac{111}{10} < 0$, so it is strictly decreasing and vanishes at most once. It does vanish once: $d(0) = (b_{1} - b_{3})t_{13} > 0$ because every $t_{k} \ge 0$, while $d(t_{13}) \le 0$. Call the zero $t_{*}$. From $0 = d(t_{*}) \le (b_{1} - b_{3})(t_{13} - t_{*})$ and $b_{1} - b_{3} > 0$ follows $t_{*} \le t_{13}$. Next, $b_{3} = -7$ and $\Dl^{(k)}_{3} \ge 0$ give
\[
a_{3}(t_{*}) \;=\; a_{3} + 7t_{*} - \sum_{k}\Dl^{(k)}_{3}(t_{*} - t_{k})_{+} \;\le\; a_{3} + 7t_{*} \;\le\; a_{3} + 7t_{13}.
\]
Equality $t_{*} = t_{13}$ holds if and only if $d(t_{13}) = 0$, that is $\sum_{k}(\Dl^{(k)}_{1} - \Dl^{(k)}_{3})(t_{13} - t_{k})_{+} = 0$; the terms are nonnegative, and $(t_{13} - t_{k})_{+} > 0$ exactly for the thresholds below the crossing, so this says $\Dl^{(k)}_{1} = \Dl^{(k)}_{3}$ at each of them. If moreover $\Dl^{(k)}_{1} = \Dl^{(k)}_{3} = 0$ at every threshold below the crossing, then $t_{*} = t_{13}$ and the sum in $a_{3}(t_{13})$ is empty of nonzero terms, so $a_{3}(t_{*}) = a_{3} + 7t_{13}$: both bounds are attained. The wino row of Proposition \claimref{standard-model.table-exact}, $\Dl = (0, \tfrac43, 0)$, is such content. A meeting of the three couplings is in particular a point with $a_{1} = a_{3}$, so it occurs at $t_{*}$.

\emph{(ii).} At a meeting, $M_{U} = M_{Z}e^{2\pi t_{*}} \le M_{Z}e^{2\pi t_{13}} = M_{13}$ and $\alpha_{U}^{-1} = a_{3}(t_{*}) \le a_{3} + 7t_{13}$ by (i). For the scale, the one transcendental step: with the record value $M_{Z} = 91.1876$ GeV $< 91.19$ GeV and $\pi < \tfrac{22}{7}$, $e^{2\pi t_{13}} < e^{44t_{13}/7}$, and a Taylor majorant of the exponential evaluated in exact rationals gives $M_{13} < 2.5 \times 10^{14}$ GeV (the central value is $2.46 \times 10^{14}$ GeV; with $\pi > \tfrac{223}{71}$ and $M_{Z} > 91.18$ GeV the matching minorant gives $M_{13} > 2.4 \times 10^{14}$ GeV). The proton figure $P := \tfrac{\alpha_{U}^{-1}}{45}(M_{U}/10^{15}\,\mathrm{GeV})^{2}$ of Fact \claimref{standard-model.imported-proton-condition} is nonpositive if $\alpha_{U}^{-1} \le 0$, and otherwise increasing in both $\alpha_{U}^{-1}$ and $M_{U}$, so in every case
\[
P \;\le\; \frac{a_{3} + 7t_{13}}{45}\Bigl(\frac{M_{13}}{10^{15}\,\mathrm{GeV}}\Bigr)^{2} \;<\; \frac{165219412}{4093125 \cdot 45} \cdot \frac{1}{16} \;<\; 0.057 \;<\; \frac{1}{10}.
\]
The condition demands $P > 5.2\sqrt{15}$, and $5.2\sqrt{15} > 20$ because $(5.2)^{2} \cdot 15 = 405.6 > 400$. So no content satisfying the hypothesis, at any number of thresholds, meets the condition, and the figure misses by a factor $5.2\sqrt{15}/P > 20 \cdot 10 = 200$. The lifetime of the Fact is proportional to $P^{2}$, so the lifetime misses by a factor above $(20/0.057)^{2} > 1.2 \times 10^{5}$; at the central $M_{13}$ the figure is $P = 0.054$ and the factor $1.4 \times 10^{5}$, about $10^{5}$.

\emph{(iii).} The three admissible rows of Proposition \claimref{standard-model.table-exact}(d) are one-threshold spectra with thresholds above $M_{Z}$, $t_{I} \ge 0$, and $\Dl = (0, \tfrac43, 0)$, $(\tfrac25, 2, 0)$, $(\tfrac52, \tfrac{25}{6}, 2)$ satisfy $\Dl_{1} \ge \Dl_{3} \ge 0$ ($0 \ge 0$, $\tfrac25 \ge 0$, $\tfrac52 \ge 2$). So (ii) applies to each, and none meets the proton condition. This uses only the running of the Fact, the record values, $b_{3} < 0$ and the rows' $\Dl$, and so is an argument for clause (e) of that proposition independent of the candidate-by-candidate evaluation it cites. The rows' meeting scales of $1.5$ to $2.5 \times 10^{14}$ GeV and Theorem \claimref{standard-model.one-threshold-exact-criterion}(i), which places each row's unique meeting, are consistent with $M_{U} \le M_{13}$, the wino row sitting on the bound as (i) says.

\emph{The closing remark.} By Fact \claimref{standard-model.imported-coefficient-sum-indices}, Weyl fermions and complex scalars added at a threshold contribute $\Dl_{i} = \tfrac23\sum_{\mathrm{Weyl}}T_{i}(\rho) + \tfrac13\sum_{\mathrm{scalars}}T_{i}(\rho)$, additive over the fields, with every index $T_{i}(\rho) \ge 0$; hence physical content has $\Dl_{3} \ge 0$ (and $\Dl_{1}, \Dl_{2} \ge 0$). A colour singlet has $T_{3} = 0$, so $\Dl_{1} \ge 0 = \Dl_{3}$. A complete multiplet with one threshold has $\Dl = c(1,1,1)$ (the first row of Proposition \claimref{standard-model.table-exact}), so $\Dl_{1} = \Dl_{3}$. A colour triplet in an $SU(2)$ representation of dimension $d$ with hypercharge $Y$, counted with weight $\kappa = \tfrac23$ (Weyl) or $\tfrac13$ (scalar), has $T_{3} = \tfrac12 d$ and $T_{1} = \tfrac35 Y^{2} \cdot 3d$, so
\[
\Dl_{1} - \Dl_{3} \;=\; \kappa d\bigl(\tfrac95 Y^{2} - \tfrac12\bigr) \;\ge\; 0 \iff Y^{2} \ge \tfrac{5}{18} \iff |Y| \ge \sqrt{5/18}.
\]
The $Q$-type triplet ($Y = \tfrac16$, $Y^{2} = \tfrac{1}{36}$) and the $d^{c}$-type triplet ($Y = \tfrac13$, $Y^{2} = \tfrac19$) fall below $\tfrac{5}{18}$, and the gluino, a Weyl adjoint of $SU(3)$ with $T_{3} = 3$, has $\Dl = (0, 0, 2)$, so $\Dl_{1} - \Dl_{3} = -2$. Finally, a one-loop spectrum that meets the proton condition has $\Dl^{(k)}_{3} \ge 0$ at every threshold, being physical content, so by (ii) it cannot have $\Dl^{(k)}_{1} \ge \Dl^{(k)}_{3}$ at every threshold: some threshold carries content with $\Dl_{3} > \Dl_{1}$.

The check recomputes $t_{13}$, $a_{3} + 7t_{13}$ and the bounds on $M_{13}$, $P$ and the two miss factors in exact rationals; verifies (i) on multi-threshold spectra by locating the crossing exactly; confirms the three admissible rows satisfy the hypothesis; and rebuilds the triplet criterion and the gluino's $\Dl$ from the index formula.
